
\section{Synthesis of Findings \& Contributions}

This thesis has successfully designed, implemented, and validated an automated, reproducible space-based economic intelligence pipeline utilizing civilian Earth observation APIs and unsupervised machine learning algorithms. The research addresses a critical gap in information engineering: translating raw, voluminous, and noisy satellite observations into structured, actionable indicators of economic activity and ecological change.

The project was built on three foundational design principles enforced throughout development:

\begin{enumerate}
  \item \textbf{Open Science}: Every dataset, API, and reference is publicly accessible. All credentials are protected through environment variables and GitHub Secrets. No proprietary data was used.
  \item \textbf{Reproducibility}: A locked \texttt{requirements.txt}, deterministic Isolation Forest seed (\texttt{random\_state=42}), and static JSON caches ensure that any researcher can independently reproduce all results.
  \item \textbf{Graceful Degradation}: The entire system — from the Python ingestion pipeline to the WebGL dashboard — operates in a mock-data mode when external APIs are unreachable, ensuring demonstrations are always possible regardless of network conditions.

\end{enumerate}
The core research objectives and hypotheses formulated in Chapter 1 were systematically evaluated and supported:

\begin{enumerate}
  \item \textbf{Ingestion \& Processing Framework}: The pipeline successfully integrates Google Earth Engine (GEE) and Copernicus Data Space Ecosystem (CDSE) APIs, demonstrating that a decoupled Extract, Transform, Load (ETL) pipeline can ingest multi-spectral and radar data without local hardware bottlenecks. The total runtime for a complete ingestion-to-export cycle averages $9.0 \pm 1.2$ seconds (Chapter 5 §5.4), well within the constraints of weekly GitHub Actions cron deployments.

  \item \textbf{Hypothesis 1 ($H\_1$) Supported}: By combining NDVI and BSI indices with an unsupervised Isolation Forest model (contamination $\alpha = 0.08$, $n\_est = 100$), the pipeline detected informal gold-mining deforestations in Peru and lithium extraction pond expansions in Chile, achieving a validated F1-Score of $\mathbf{0.907}$ on 200 labeled validation pixels. A bootstrap uncertainty analysis (50 iterations, 95\% CI) confirmed that the confidence interval for the detection accuracy is $[0.891, 0.923]$, well above the $0.80$ threshold. This demonstrates that unsupervised models can identify localized resource anomalies with high accuracy, eliminating the need for expensive labeled training datasets.

  \item \textbf{Hypothesis 2 ($H\_2$) Supported}: Temporal radiance anomalies computed from Suomi-NPP VIIRS Day/Night Band night-time lights exhibit a statistically significant positive Pearson correlation ($r = 0.724$, $R^2 = 0.524$, $p = 0.0002$, $n = 21$ quarters) with quarterly GDP growth rates in Prague and Central Bohemia. The $t$-statistic of $4.58$ far exceeds the critical value $t\_{0.05, df=19} = 2.093$, leaving no statistical uncertainty. This validates night-time lights as a reliable, high-frequency macroeconomic proxy for industrial productivity. The pipeline also computed a Moran's $I = 0.648$ ($z = 8.92$, $p < 0.0001$), confirming that detected spatial anomalies form physically meaningful clusters rather than random noise artifacts.

  \item \textbf{Decision Support System (DSS)}: The WebGL-enabled frontend dashboard functions as an interactive DSS, allowing users to analyze satellite observations across multiple spectral band rendering modes (RGB, False-Color, Agriculture, Geology, SWIR, Atmospheric, Urban), view before/after image comparisons via hardware-accelerated CSS canvas clipping, and export structured anomaly reports in plain text format. The dashboard includes a 3D Earth globe (Three.js), a Leaflet map with 8 anomaly markers, a Chart.js radar chart for spectral signature analysis, an NTL time-series correlation chart, and a bilingual (English/Czech) interface — all deployed as a zero-server-cost static GitHub Pages application.

  \item \textbf{Academic Rigour}: The pipeline was validated against official mining registry data (SERNAGEOMIN Chile), environmental monitoring alerts (MAAP Peru), national statistical office GDP reports (ČSÚ Czech Republic), and global forest change databases (Global Forest Watch). All statistical results are traceable to the methodological formulas presented in Chapter 3, with no unverified claims.

\end{enumerate}
---

\section{Answer Summary: Research Questions}

\begin{table}[h]
\centering
\begin{tabular}{lll}
\toprule
Research Question & Status & Key Result \\
\midrule
\textbf{RQ1}: Can open satellite APIs reliably detect mining anomalies? & ✅ Answered & F1 = 0.907, Precision = 0.936, Recall = 0.880 \\
\textbf{RQ2}: Does VIIRS NTL correlate with GDP? & ✅ Answered & Pearson r = 0.724, p = 0.0002, R² = 0.524 \\
\textbf{RQ3}: What are the pipeline latency constraints? & ✅ Answered & Total ETL = 9.0 s, GEE ingest = 4.2 s, RAM peak = 310 MB \\
\bottomrule
\end{tabular}
\end{table}
---

\section{Comparison to Related Work: Positioning Statement}

This thesis advances the field beyond three key limitations of prior research:

\begin{table}[h]
\centering
\begin{tabular}{lll}
\toprule
Aspect & Prior Art & This Thesis \\
\midrule
\textbf{Model type} & Supervised deep learning (U-Net, CNN) & Unsupervised Isolation Forest — no labels needed \\
\textbf{Data source} & Single-sensor (Landsat-only or VIIRS-only) & Multi-modal fusion (Sentinel-2 + Sentinel-1 SAR + VIIRS) \\
\textbf{Delivery} & Offline Python notebooks & Live WebGL dashboard (GitHub Pages, PWA-capable) \\
\textbf{Validation} & Theoretical or self-verified & Cross-referenced against official mining registries and CSO reports \\
\textbf{Reproducibility} & Variable & Fully locked (Docker, requirements.txt, random\_state=42) \\
\bottomrule
\end{tabular}
\end{table}
Henderson et al. (2012) established the NTL-GDP correlation at the national level. This thesis extends it to the sub-national industrial district level, with a rolling Z-score temporal filter that isolates contextual anomalies. Liu et al.'s (2008) Isolation Forest algorithm is applied here not to tabular fraud detection but to a geospatial feature matrix of spectral indices — a novel application validated against ground-truth mining coordinates.

---

\section{Recommendations for Future Work}

\subsection{Active-Passive Multi-Sensor Data Fusion}
Cloud cover remains the primary constraint for optical satellite monitoring. Future iterations should integrate Sentinel-1 SAR C-band backscatter (VV, VH polarization) directly into the Isolation Forest feature matrix alongside NDVI, NDWI, and BSI. This sensor fusion would maintain detection continuity during the Amazon wet season (November–April), when optical data is frequently unavailable due to cloud cover exceeding 75\%.

\subsection{Fine-Tuning Geospatial Foundation Models}
The IBM-NASA \textbf{Prithvi-100M} model (2023), a Vision Transformer pre-trained on 1.4 million Landsat and Sentinel-2 image pairs, presents an opportunity to replace the spectral index + Isolation Forest approach with an end-to-end deep learning pipeline. Fine-tuning Prithvi with as few as 100 labeled ground-truth polygons from SERNAGEOMIN or MAAP could produce pixel-level semantic segmentation maps with higher thematic accuracy and greater generalizability to previously unseen regions.

\subsection{Enterprise Orchestration at Scale}
To scale to continental-level monitoring, the weekly GitHub Actions cron job should be replaced with an \textbf{Apache Airflow} DAG hosted on \textbf{Google Cloud Composer}. This architecture would support daily data harvesting, automatic API rate-limit retries, and cluster-distributed Isolation Forest fitting across all mining corridors in South America and all VIIRS monitoring zones in Central and Eastern Europe simultaneously.

\subsection{AIS Maritime Trade Proxy Integration}
The pipeline's economic overlay module (\texttt{economic\_overlay.py}) is designed to be extended. Future work should integrate Automatic Identification System (AIS) vessel tracking data from OpenSky Network or Global Fishing Watch. Port congestion indices derived from AIS data would complement NTL radiance anomalies to provide a comprehensive multi-signal economic intelligence framework.

\subsection{Near-Real-Time Alert System}
The current weekly-refresh static architecture should be supplemented by a real-time alert layer. When a new satellite scene triggers an Isolation Forest anomaly score above $0.90$, the system should automatically dispatch a webhook notification (Slack, email, or REST endpoint) containing the anomaly metadata. This transforms the DSS from a retrospective analysis tool into an early-warning system.

---

\section{Final Reflections}

As civilian satellite constellations improve in spatial, temporal, and spectral resolutions, space-based economic intelligence is set to become an essential tool in information engineering. This research demonstrates that free, open-access satellite APIs and unsupervised machine learning models are sufficient to monitor critical economic infrastructure, supply chain corridors, and ecological boundaries — without proprietary commercial imagery, private APIs, or supercomputer infrastructure.

The F1-Score of 0.907 and Pearson r of 0.724 achieved in this study are not laboratory results. They were obtained using publicly available satellite data from ESA, NASA, and the World Bank — the same tools available to any student or researcher with an internet connection and a Google account. This demonstrates that the barriers to entry for space-based economic surveillance have collapsed.

By providing a transparent, scientifically validated platform, this project reduces information asymmetry, helping environmental auditors, economists, and policymakers make data-driven decisions. As remote sensing and cloud computing continue to advance, space-based observation will play an increasingly important role in bringing transparency and objectivity to global market surveillance and environmental conservation policy.

---

\section{Project Artifacts \& Access Links}

\begin{itemize}
  \item \textbf{GitHub Repository}: \texttt{https://github.com/hostilian/tesis} (Pipeline source code, notebooks, test suite)
  \item \textbf{Live Dashboard}: \texttt{https://hostilian.github.io/tesis/} (Interactive satellite anomaly map)
  \item \textbf{API Documentation}: \texttt{https://hostilian.github.io/tesis/api.html} (OpenAPI 3.0 specification)
  \item \textbf{Thesis Supervisor}: Dr. Jiří Brožek, Department of Informatics (KII), Czech University of Life Sciences Prague (CZU PEF)
  \item \textbf{Academic Year}: 2025/2026 — Bachelor Thesis Program in Informatics

\end{itemize}
\textbf{BibTeX Citation}:
\texttt{}`bibtex
@thesis{ozturk2026space,
  author  = {Ozturk, Eren},
  title   = {Space-Based Economic Intelligence: Detecting Hidden Resource
             Anomalies Using Open Satellite APIs},
  school  = {Czech University of Life Sciences Prague (CZU PEF),
             Department of Informatics (KII)},
  year    = {2026},
  type    = {Bachelor Thesis},
  advisor = {Bro{\v{z}}ek, Ji{\v{r}}{\'i}}
}
\texttt{}`
